\chapter{Introduction}\label{chap:intro}

% -----------------------------------------------------------------------------
% Citation placeholders follow the [Author, Year] style.
% Replace with \cite{bibkey} once BibTeX keys are defined in bibliography.tex.
% -----------------------------------------------------------------------------

\section{Motivation}\label{sec:motivation}

In a small classroom, a good teacher rarely relies on test scores alone. A furrowed brow during a derivation, a distant gaze halfway through an explanation, a brief smile after a difficult problem, and similar cues quietly tell the teacher when to slow down, rephrase, or move on. Decades of research in affective neuroscience and educational psychology confirm that this teacherly intuition has deep roots. Emotion is not a byproduct of thinking but an integral part of it, and learning, as it unfolds in schools and in the real world, is not a disembodied or purely rational process [Immordino-Yang and Damasio, 2007]. Emotional states modulate attention, encoding and memory consolidation, so that what a student feels while studying directly shapes what they later remember and how they will reason about it [Tyng et al., 2017]. In tutoring contexts specifically, a distinctive family of academic emotions, namely engagement, confusion, frustration, and boredom, emerges repeatedly, each carrying pedagogical meaning of its own [D'Mello and Graesser, 2012]. Pekrun's Control-Value Theory explains \emph{why} these emotions matter: learners' appraisals of how much control they have over a task and how much they value it arouse achievement emotions that in turn feed back into engagement, motivation, and ultimately performance [Pekrun, 2006]. A learner's face, therefore, is not decorative information; it is a window into the very processes that education is trying to shape.

However, the contextual cues that a skilled teacher reads almost unconsciously largely disappear when learning moves online or to large-scale digital platforms. Modern e-learning and intelligent tutoring environments have expanded rapidly, especially after the pandemic, and are now a central channel through which millions of students study [Strielkowski et al., 2025]. Yet these systems typically treat learners as receivers of information, with little concern for their moment-to-moment emotional or cognitive condition, and adverse states such as boredom, anxiety, or confusion can therefore persist silently and significantly hinder learning performance [Gong, 2026]. Studies of online and STEM learning repeatedly show that emotions shape engagement, motivation, and satisfaction, and that ignoring them leaves platforms unable to respond when students disengage [Anwar et al., 2023]. The effect is compounded by the fact that supervision is thin in online settings: when nobody notices that attention is drifting, sustained inattention tends to be one of the first things that breaks down, eroding the very benefits that the medium was meant to deliver [Gupta et al., 2023; Aly et al., 2023]. The picture that emerges is a paradox: precisely in the environments where personalization is most needed, the richest source of human feedback, the learner's face, is the one that digital systems most systematically ignore.

\section{Problem Statement}\label{sec:problem}

A widely proposed response to this paradox is adaptive learning: systems that tailor the sequence, difficulty, feedback, and pacing of material to each individual learner rather than delivering a fixed path [Gligorea et al., 2023; Strielkowski et al., 2025]. Recent intelligent tutoring systems (ITS) have taken up this promise with increasingly sophisticated tools. Reinforcement-learning approaches have been used to adjust task difficulty and instructional complexity in simulated learners [Axak et al., 2025; Deshmukh and Sen, 2025], and have been combined with knowledge tracing to optimise personalised learning paths across students' evolving states [Fu, 2025]. Comparative studies show that such policies can meaningfully improve task selection over rule-based baselines [Azoulay et al., 2021], and recent work couples knowledge tracing with reinforcement-learning-based problem selection to produce measurable learning gains in short human studies [Guzman and Cruz-Mercado, 2025]. Even outside the reinforcement-learning paradigm, biologically inspired memory architectures have been used to recommend the next exercise within each learner's Zone of Proximal Development [Kuling and Zitnik, 2025]. Despite this progress, two patterns emerge clearly from recent surveys of the field. First, a surprising number of adaptive systems are described in the literature but remain experimental, with many not validated in authentic classroom settings and often disconnected from the concrete problems students face [Kabudi et al., 2021; Zerkouk et al., 2025]. Second, and more importantly for this thesis, most adaptive systems still operate only on cognitive-behavioural traces, such as answers, response times, clicks, and hint requests, and treat the learner as an affect-free agent [Axak et al., 2025; Fu, 2025; Azoulay et al., 2021]. Broader comparative work echoes that behavioural traces predominate when signals are chosen for the policy loop [Guzman and Cruz-Mercado, 2025; Weitekamp et al., 2025]. Urbaite warns that a persistent risk in such systems is collapsing the multidimensional construct of engagement into ``easy-to-log proxies'' such as clicks and time-on-page [Urbaite, 2026]. Meng and Yang similarly observe that existing RL-driven tutoring solutions often overlook engagement factors altogether and, when they do attempt to model them, rely primarily on textual interaction patterns rather than on perceptual signals of the student [Meng and Yang, 2025]. Authors of recent RL-based tutors explicitly flag the integration of affective computing as open future work, acknowledging that their policies are, for now, emotionally blind [Axak et al., 2025; Guzman and Cruz-Mercado, 2025].

Facial Expression Recognition (FER) offers a practical way to recover part of the signal that adaptive systems currently ignore. Over the last decade, deep learning has substantially advanced FER accuracy, with convolutional networks, vision transformers, and hybrid temporal models all contributing to stronger recognition in unconstrained settings [Lek and Teo, 2023; Liu et al., 2026; Aly, 2025]. Attention-heavy variants further tighten performance in challenging conditions [Aly and Alotaibi, 2025]. In educational contexts, researchers have moved beyond Ekman's six basic emotions, which rarely occur in actual study sessions, towards academic affect dimensions such as engagement, boredom, confusion, and frustration, the very states singled out by earlier discourse-analytic work on tutoring [Leong, 2020; D'Mello and Graesser, 2012]. The DAiSEE dataset has become a common anchor for this line of research, and a growing body of deep-learning models, ranging from facial-embedding and landmark approaches to more recent dynamic FER models with attention and temporal mechanisms, has been trained on it [Leong, 2020; Asaju and Vadapalli, 2022; Mehta et al., 2022]. Related architectures and training regimes continue to expand in the latest DAiSEE-based studies [Komaravalli and Janet, 2023; Liu et al., 2026; Rao et al., 2021]. Parallel efforts in online-learning contexts use modern CNN and transformer architectures to monitor students in live classes [Aly, 2025; Aly et al., 2023]. Additional classroom-oriented deployments extend the same design space [Aly and Alotaibi, 2025; Gupta et al., 2023]. Systematic reviews of academic FER confirm that deep learning now dominates this subfield and that real-time deployment is becoming realistic, while also noting ongoing challenges from pose, illumination, occlusion, and cross-cultural variability [Lek and Teo, 2023; Fernández-Herrero, 2024]. Overall, the evidence that facial signals carry meaningful information about engagement and cognitive state in online learners has become difficult to dispute [Mehta et al., 2022; Gupta et al., 2023; Woolf et al., 2023].

Some recent work does connect affect to adaptation, including reinforcement-learning formulations that adjust difficulty or pacing from inferred emotional state. In practice, however, many of these systems lean on multimodal or physiological channels, such as heart rate, EEG, respiration, text, voice, or tone [Gong, 2026; Govea et al., 2024]. Such setups are often heavy in practice, since signals must be acquired, synchronised, and processed alongside the teaching interaction, which typically implies specialised hardware and careful integration before the approach can be deployed at scale in ordinary online courses [Gutierrez et al., 2025; Kong et al., 2025; Nandimandalam, 2025]. A second limitation, within this same emotion-aware adaptation literature, is that affect is not always carried through to the decision process that selects the next task; many systems stop at detection or dashboard-style reporting, so the intended loop from detecting emotion, to using it in the adaptation decision, to adjusting the task or its difficulty is rarely closed [Fern\'andez-Herrero, 2024; Septiana et al., 2024; Vistorte et al., 2024]. Using facial expression alone offers a lighter alternative that is less intrusive than body-worn sensors and realistic to collect with commodity webcams for large cohorts. Prior reinforcement-learning work that uses facial affect as feedback for adaptation shows that such a loop is feasible, but evaluation is often simulation-based and the affective signal may be coarse or averaged across whole tasks rather than supplied as a continuous input at decision time [P\'erez et al., 2024]. This thesis therefore adopts an FER-only, webcam-based route from the student's facial signal to the next task decision.

While facial expression recognition in education and adaptive tutoring have both matured, they have largely matured in parallel. Most FER-for-learning systems stop at recognition: they train a classifier on DAiSEE or a similar dataset, report per-class accuracy, and at best propose an analytics dashboard or teacher-facing feedback channel [Leong, 2020; Rao et al., 2021; Mehta et al., 2022]. Studies that span DAiSEE-trained models to online prototypes typically remain at the recognition stage for design reasons [Komaravalli and Janet, 2023; Liu et al., 2026; Aly, 2025]. Further deployments discussed in the same literature often mirror that scope [Aly et al., 2023]. Where these systems do discuss adaptation, it is typically aspirational, leaving the pedagogical response either to a human instructor or to unspecified future work [Asaju and Vadapalli, 2022; Aly and Alotaibi, 2025; Prabakaran et al., 2025]. In parallel, adaptive tutoring systems that do implement closed-loop decisions rarely use facial affect on its own; as noted above, the face, when included, tends to be one channel among several, with non-facial signals typically carrying most of the weight. Scoping reviews make the same observation at a field level. Affective intelligent tutoring systems, it has been noted, ``recognize more than they regulate'', still leaning on single-modality detection and rarely closing the loop with real-time interventions [Yuvaraj et al., 2025]. There is, in short, a clear lack of complete FER-to-adaptation pipelines in which facial expressions drive a concrete task-sequencing or difficulty-adjustment policy, together with an evaluation that actually isolates the contribution of the emotional signal. This thesis is motivated by that gap.

\section{Objectives}\label{sec:objectives}

This thesis investigates whether facial-expression-derived academic emotion can be used as an additional signal in an adaptive learning system to inform decisions about the next task and its difficulty. The central hypothesis is that augmenting a sequencing policy with an FER-derived affective input leads to more pedagogically sensitive adaptation than a policy that relies on cognitive-behavioural signals alone. The work pursues the following objectives:

\begin{enumerate}
  \item To develop and evaluate a focused FER module for academic emotions in an educational setting, building on architectures that have proven effective on DAiSEE and related data [Mehta et al., 2022; Liu et al., 2026; Aly and Alotaibi, 2025].
  \item To design a task-sequencing mechanism that treats FER output as an explicit input alongside learner performance, so that affect is part of the adaptation policy rather than only an analytics layer.
  \item To evaluate the emotion-aware policy against an emotion-free baseline and thereby isolate the contribution of the facial signal, in line with calls for affective systems that go beyond recognition to act on what they perceive [Fern\'andez-Herrero, 2024; Yuvaraj et al., 2025; Lek and Teo, 2023].
\end{enumerate}

\section{Thesis Outline}\label{sec:outline}

The remainder of this thesis is organised as follows. A background chapter reviews the theoretical foundations of emotion in learning, the state of the art in facial expression recognition, and the methods used in adaptive learning systems. The methodology chapter presents the design of the FER module, the adaptation mechanism, and their integration into a single pipeline. The results chapter reports the experimental setup and quantitative findings, first for the FER component in isolation and then for the adaptive policy with and without the emotional signal. The discussion and conclusion chapters interpret the findings in relation to the gap identified above, discuss limitations and ethical considerations, and outline directions for future work.
